\documentclass[a4paper,11pt]{article}

\usepackage[T1]{fontenc}
\usepackage[utf8]{inputenc}
\usepackage{geometry}
\usepackage{array}
\usepackage{float}
\usepackage{longtable}
\usepackage{listings}
\usepackage{natbib}
\usepackage{siunitx}
\usepackage[hidelinks]{hyperref}

\hypersetup{
  pdftitle={Online Resource 1: Supplementary Material for LLM-Based SQL Antipattern Detection in jOOQ Code: An Occurrence-Level Evaluation and Repository Study},
  pdfauthor={Kristo Isberg; Erki Eessaar},
  pdfsubject={Supplementary information for Empirical Software Engineering},
  pdfkeywords={SQL antipatterns, jOOQ, large language models, reproducibility, supplementary material}
}

\geometry{margin=25mm}
\sisetup{mode=text}

\title{Online Resource 1: Supplementary Material for\\
\emph{LLM-Based SQL Antipattern Detection in jOOQ Code:\\
An Occurrence-Level Evaluation and Repository Study}}
\author{Kristo Isberg \and Erki Eessaar}
\date{}

\begin{document}
\maketitle

\begin{center}
\emph{Empirical Software Engineering}\\[0.5em]
Kristo Isberg, School of Information Technologies, Tallinn University of Technology, Tallinn, Estonia\\
Erki Eessaar, Department of Software Science, School of Information Technologies, Tallinn University of Technology, Tallinn, Estonia\\
Corresponding author: Kristo Isberg, \href{mailto:krisbe@taltech.ee}{krisbe@taltech.ee}
\end{center}

This supplement contains supporting detail and indexes the larger frozen replication resources. The full study artefacts are preserved at commit \href{https://github.com/kristoisberg/masters-thesis/tree/d9b35e398a6deb544f913bdbc0b211ab38474a44}{\texttt{d9b35e3}}. Detector source is preserved separately at commit \href{https://github.com/kristoisberg/jooq-antipattern-detector/tree/cf82fe56acb728f076df67279ff4a78a138996f3}{\texttt{cf82fe5}}. Article analyses and generated corpus-normalisation tables are preserved at commit \href{https://github.com/kristoisberg/thesis-paper/tree/bb3bf6056cc7bbda29bb0fb5f0b720e7f9d06028/analysis}{\texttt{bb3bf60}}.

\section{Resource index}

\begin{description}
    \item[Repository mining.] The frozen repository contains the \href{https://github.com/kristoisberg/masters-thesis/blob/d9b35e398a6deb544f913bdbc0b211ab38474a44/thesis/appendices/appendix-github-search-terms.tex}{GitHub search strings}, \href{https://github.com/kristoisberg/masters-thesis/blob/d9b35e398a6deb544f913bdbc0b211ab38474a44/thesis/appendices/appendix-omitted-projects.tex}{omitted-project list}, and mining notebooks under \href{https://github.com/kristoisberg/masters-thesis/tree/d9b35e398a6deb544f913bdbc0b211ab38474a44/scripts}{\texttt{scripts/}}.
    \item[Annotation materials.] The \href{https://github.com/kristoisberg/masters-thesis/blob/d9b35e398a6deb544f913bdbc0b211ab38474a44/thesis/appendices/appendix-annotated-antipatterns.tex}{full 19-class catalogue}, \href{https://github.com/kristoisberg/masters-thesis/tree/d9b35e398a6deb544f913bdbc0b211ab38474a44/diagrams/decision}{decision trees}, and reference datasets are preserved in the frozen repository.
    \item[Detector configuration.] All query and schema \href{https://github.com/kristoisberg/masters-thesis/tree/d9b35e398a6deb544f913bdbc0b211ab38474a44/prompts}{prompts}, model-selection detail, per-class configuration results, and configuration options are preserved with the study artefacts.
    \item[Evaluation and corpus analysis.] The frozen notebooks contain the localisation and classification diagnostics, detector-informed sensitivity analysis, and co-detection matrices and heatmaps. The ordered source-fragment categorisation procedure and stored counts are in \href{https://github.com/kristoisberg/masters-thesis/blob/d9b35e398a6deb544f913bdbc0b211ab38474a44/scripts/21-find-frequent-offenders.ipynb}{\texttt{scripts/21-find-frequent-offenders.ipynb}}. The article repository contains the \href{https://github.com/kristoisberg/thesis-paper/blob/bb3bf6056cc7bbda29bb0fb5f0b720e7f9d06028/analysis/corpus_normalization.py}{corpus-normalisation script}, \href{https://github.com/kristoisberg/thesis-paper/blob/bb3bf6056cc7bbda29bb0fb5f0b720e7f9d06028/analysis/corpus_phase2_class_summary.csv}{class-level breadth and repetition table}, \href{https://github.com/kristoisberg/thesis-paper/blob/bb3bf6056cc7bbda29bb0fb5f0b720e7f9d06028/analysis/corpus_phase3_density_summary.csv}{repository-density table}, \href{https://github.com/kristoisberg/thesis-paper/blob/bb3bf6056cc7bbda29bb0fb5f0b720e7f9d06028/analysis/corpus_phase3_concentration.csv}{concentration decomposition}, and \href{https://github.com/kristoisberg/thesis-paper/blob/bb3bf6056cc7bbda29bb0fb5f0b720e7f9d06028/analysis/corpus_phase4_exact_duplicate_sensitivity.csv}{exact-content sensitivity table}. Sections~\ref{sec:splitSupport} and~\ref{sec:iouSensitivity} provide stand-alone copies of two compact diagnostics also reported in the article.
    \item[Implementation.] Architecture, workflow, output, and command-line configuration details accompany the frozen \href{https://github.com/kristoisberg/jooq-antipattern-detector/tree/cf82fe56acb728f076df67279ff4a78a138996f3}{detector source} and the original thesis appendices.
\end{description}

\section{Reproducibility record}

This record distinguishes preserved evidence from information that the original executions did not capture. No DOI-backed replication deposit existed when this record was prepared. The immutable commits below identify the surviving material. A repository DOI remains required.

\clearpage
\subsection{Snapshots and preserved outputs}

Table~\ref{tab:reproSnapshots} maps each immutable snapshot to the evidence it preserves and the limits of that evidence.

\begin{longtable}{>{\raggedright\arraybackslash}p{0.20\textwidth}>{\raggedright\arraybackslash}p{0.24\textwidth}>{\raggedright\arraybackslash}p{0.46\textwidth}}
    \caption{Immutable source snapshots and preserved outputs. The identifier and date locate each snapshot; the final column states the evidence retained and its preservation boundary.}\label{tab:reproSnapshots}\\
    \hline
    \textbf{Resource} & \textbf{Identifier and archive date} & \textbf{Preserved evidence and boundary} \\ \hline
    \endfirsthead
    \caption[]{Immutable source snapshots and preserved outputs (continued).}\\
    \hline
    \textbf{Resource} & \textbf{Identifier and archive date} & \textbf{Preserved evidence and boundary} \\ \hline
    \endhead
    \hline
    \endfoot
    Study artefacts & \href{https://github.com/kristoisberg/masters-thesis/tree/d9b35e398a6deb544f913bdbc0b211ab38474a44}{\texttt{d9b35e3}}, 24 May 2026 & \path{datasets/test-set.csv} contains the \num{523} original test references. The executed \path{scripts/14-evaluate-tool-localisation.ipynb} contains all \num{536} held-out predictions. \path{datasets/analysis-results.csv} contains all \num{15931} positive corpus flags; per-project positive outputs also survive. \\
    Released detector & \href{https://github.com/kristoisberg/jooq-antipattern-detector/tree/cf82fe56acb728f076df67279ff4a78a138996f3}{\texttt{cf82fe5}}, 26 May 2026 & Source, tests, \path{package.json}, and \path{bun.lock}. This release postdates the April executions and does not identify the binary used for those runs. \\
    Article analyses & \href{https://github.com/kristoisberg/thesis-paper/tree/bb3bf6056cc7bbda29bb0fb5f0b720e7f9d06028/analysis}{\texttt{bb3bf60}}, 12 September 2026 & \path{analysis/localisation_robustness.py}, \path{analysis/corpus_concentration.py}, \path{analysis/corpus_normalization.py}, and \path{analysis/requirements.txt} reproduce the article's localisation robustness, source-frame reconstruction, normalised corpus measures, and exact-content sensitivity from the frozen study artefacts and retained source snapshot. \\
\end{longtable}

The preserved prediction tables contain the parsed class, file, span, code fragment, and explanation fields used in the analyses. Raw API response envelopes, request identifiers, negative-file outputs, failed responses, and complete request metadata were not archived. The source snapshot for the \num{602} analysed repositories survives without Git history. The corpus manifest records a repository-relative path and SHA-256 content hash for every selected file. The detector revision used for the executions remains unavailable.

\subsection{Models, dates, settings, and retries}

The validation notebooks called OpenRouter \citep{openRouter2026homepage} through the OpenAI Python SDK \citep{openAiPython2026GitHub}. They preserve four model slugs: \path{openai/gpt-5.2} \citep{gpt52}, \path{z-ai/glm-5} \citep{DBLP:journals/corr/abs-2602-15763}, \path{anthropic/claude-opus-4.5} \citep{opus45Card}, and \path{openai/gpt-oss-120b} \citep{DBLP:journals/corr/abs-2508-10925}. Table~\ref{tab:validationRuns} gives the dashboard date and aggregate retry count for every model--prompt run. Dashboard time-zone information was not saved. The surviving notebook copies process \num{1159} files and do not reproduce the file count, costs, or runtimes in the article's originally reported validation table; the exact run outputs underlying that table are unavailable.

\begin{table}[H]
    \centering
    \small
    \caption{Validation execution date in 2026 and aggregate retries, shown as date/retries. ZS is Zero-Shot, FS is Few-Shot, CoT is Chain-of-Thought, and ToT is Tree-of-Thought.}
    \label{tab:validationRuns}
    \begin{tabular}{lrrrr}
        \hline
        \textbf{Prompt} & \textbf{GPT-5.2} & \textbf{GLM-5} & \textbf{Opus 4.5} & \textbf{gpt-oss-120B} \\
        \hline
        ZS  & 6 Mar/1  & 6 Mar/175 & 7 Mar/0 & 7 Mar/0 \\
        FS  & 7 Mar/0  & 7 Mar/9   & 7 Mar/0 & 6 Mar/1 \\
        CoT & 6 Mar/1  & 6 Mar/234 & 7 Mar/0 & 6 Mar/3 \\
        ToT & 27 Mar/2 & 27 Mar/14 & 27 Mar/9 & 27 Mar/58 \\
        \hline
    \end{tabular}
\end{table}

GPT-5.2 requests specified xhigh reasoning and temperature \num{0.0}; the archive does not show whether the backend applied the temperature value. GLM-5 used temperature \num{0.0} without a reasoning parameter, fixed Friendli as the backend, and disabled fallbacks. Claude Opus 4.5 used temperature \num{0.0} with reasoning disabled. gpt-oss-120B omitted the temperature parameter, requested high reasoning, and returned metadata with temperature \num{1.0}. Each model--prompt configuration ran once. The retry totals count empty or failed responses repeated by the notebooks.

The held-out execution used \texttt{openrouter:anthropic/claude-opus-4.5}, temperature \num{0.0}, reasoning disabled, and the Zero-Shot prompts. Its dashboard record is dated 4 April 2026 at 15:00 and contains \num{502} requests for \num{502} files. The command allowed up to \num{1000} retries per file; the equal request and file counts show no additional attempts in the dashboard total. The corpus execution used the same model and generation settings on 15 April 2026 from 08:00 through 12:00. Its command allowed two retries per file, and the dashboard records \num{18051} requests. This exceeds the reconstructed \num{17988}-file frame by \num{63} requests, consistent with up to \num{63} repeated attempts. Per-request retry histories, final failure records, backend provider identifiers, and model snapshot identifiers are unavailable.

The eight query and schema prompt files survive under \path{prompts/zero-shot/}, \path{prompts/few-shot/}, \path{prompts/chain-of-thought/}, and \path{prompts/tree-of-thought/}. The validation notebooks load these files directly. The held-out and corpus runs used prompts embedded in an unarchived binary, so byte identity with the preserved Zero-Shot files cannot be established.

\subsection{Dependencies and analysis scripts}

The original notebook archive contains 26 top-level notebooks and 16 model--prompt backups. Notebook metadata records Python \texttt{3.14.3} in 38 notebooks and \texttt{3.13.5} in four backups. The notebooks import NumPy, pandas, scikit-learn, matplotlib, seaborn, the OpenAI Python SDK, Pydantic, and \texttt{requests\_async}, but the original environment has no requirements or lock file. Exact original package versions are therefore unavailable. The released detector records Bun \texttt{1.2.21} and pins its JavaScript dependency graph in \path{bun.lock}; this later release cannot establish the original run environment.

The article's reconstruction scripts were verified with Python \texttt{3.14.7}. Their direct third-party dependencies are pinned in \path{analysis/requirements.txt}. The corpus-normalisation script uses the Python standard library. From the article repository root, the three commands are:

\begin{lstlisting}[frame=single,basicstyle=\small\ttfamily]
python analysis/localisation_robustness.py /path/to/masters-thesis \
  --bootstrap 10000 --seed 123456
python analysis/corpus_concentration.py \
  /path/to/masters-thesis/datasets/analysis-results.csv \
  --verify-frozen
python analysis/corpus_normalization.py /path/to/repositories \
  /path/to/masters-thesis/datasets/final-repositories-corrected.csv \
  /path/to/masters-thesis/datasets/analysis-results.csv \
  --output-dir analysis
\end{lstlisting}

The first command reconstructs the \num{536} predictions from the executed notebook, reproduces the primary \num{460}/\num{76}/\num{63} true-positive/false-positive/false-negative totals, evaluates four intersection-over-union thresholds, and runs the \num{10000}-iteration project bootstrap with seed \num{123456}. The second verifies the frozen corpus CSV hash and reproduces the original class-level repository-concentration results. The third scans the retained source snapshot, verifies the allowlist and frozen flags, reconstructs the \num{17988}-file frame, and writes all Phase 1--4 normalisation artefacts. Every phase has assertion-based acceptance checks and a JSON summary with methods and checksums.

\subsection{Corpus-normalisation artefacts and checksums}

The generated tables and their SHA-256 values are recorded below. The linked analysis directory also contains the Phase 1--4 JSON summaries, input hashes, selection-rule digest, methods, and acceptance checks.

\begin{longtable}{>{\raggedright\arraybackslash}p{0.50\textwidth}>{\raggedright\arraybackslash}p{0.40\textwidth}}
    \caption{Corpus-normalisation artefacts at article commit \texttt{bb3bf60}. Each hash is split across two lines for typesetting. The \href{https://github.com/kristoisberg/thesis-paper/tree/bb3bf6056cc7bbda29bb0fb5f0b720e7f9d06028/analysis}{analysis directory} contains the JSON records of methods and acceptance checks.}\label{tab:normalizationArtefacts}\\
    \hline
    \textbf{Artefact} & \textbf{SHA-256} \\ \hline
    \endfirsthead
    \caption[]{Corpus-normalisation artefacts and SHA-256 values (continued).}\\
    \hline
    \textbf{Artefact} & \textbf{SHA-256} \\ \hline
    \endhead
    \hline
    \endfoot
    \href{https://github.com/kristoisberg/thesis-paper/blob/bb3bf6056cc7bbda29bb0fb5f0b720e7f9d06028/analysis/corpus_source_manifest.csv}{\path{analysis/corpus_source_manifest.csv}} & \shortstack[l]{\texttt{8200f8990215912efba90a4f5ed431d5}\\\texttt{2c79d7b20e11ce496457f34829c8a214}} \\
    \href{https://github.com/kristoisberg/thesis-paper/blob/bb3bf6056cc7bbda29bb0fb5f0b720e7f9d06028/analysis/corpus_flag_alignment.csv}{\path{analysis/corpus_flag_alignment.csv}} & \shortstack[l]{\texttt{b035fd06e7c829da8f995770e769c7d3}\\\texttt{491d0d32794837ddc2191f960f78c728}} \\
    \href{https://github.com/kristoisberg/thesis-paper/blob/bb3bf6056cc7bbda29bb0fb5f0b720e7f9d06028/analysis/corpus_phase2_by_repository.csv}{\path{analysis/corpus_phase2_by_repository.csv}} & \shortstack[l]{\texttt{a88de1ee4719d1f8290d1418940a1115}\\\texttt{1fb7cbe443023895002141493e4352a2}} \\
    \href{https://github.com/kristoisberg/thesis-paper/blob/bb3bf6056cc7bbda29bb0fb5f0b720e7f9d06028/analysis/corpus_phase2_class_summary.csv}{\path{analysis/corpus_phase2_class_summary.csv}} & \shortstack[l]{\texttt{7a91098968dc2b20afd86edf3cc505c2}\\\texttt{cd21fcd220818e5888a7c144639e451b}} \\
    \href{https://github.com/kristoisberg/thesis-paper/blob/bb3bf6056cc7bbda29bb0fb5f0b720e7f9d06028/analysis/corpus_phase3_density_summary.csv}{\path{analysis/corpus_phase3_density_summary.csv}} & \shortstack[l]{\texttt{598681f34f558054c2b6cb9f92daac65}\\\texttt{6e8dfb2f76a5e85069ad2a9866236334}} \\
    \href{https://github.com/kristoisberg/thesis-paper/blob/bb3bf6056cc7bbda29bb0fb5f0b720e7f9d06028/analysis/corpus_phase3_concentration.csv}{\path{analysis/corpus_phase3_concentration.csv}} & \shortstack[l]{\texttt{7a5c900e7d22e3e638a7d12960203ee2}\\\texttt{39902e175345f1d0a3d5a41ecb1fd2c9}} \\
    \href{https://github.com/kristoisberg/thesis-paper/blob/bb3bf6056cc7bbda29bb0fb5f0b720e7f9d06028/analysis/corpus_phase4_exact_duplicate_sensitivity.csv}{\path{analysis/corpus_phase4_exact_duplicate_sensitivity.csv}} & \shortstack[l]{\texttt{15c16d4e23a18faa68dbe3fe9c5f931d}\\\texttt{7620f0194650f04c226e351fb985355d}} \\
\end{longtable}

Table~\ref{tab:normalizationArtefacts} provides direct links to the complete source, alignment, class, repository, density, concentration, and exact-content tables.

\section{Equivalent SQL and jOOQ source representations}

The following fragments express the same Implicit Columns occurrence in SQL and jOOQ \citep{jooq2025sql}: both select every column from the \texttt{employee} table.

\paragraph{SQL}
\begin{lstlisting}[frame=single,basicstyle=\small\ttfamily]
SELECT * FROM employee WHERE name = 'Jack'
\end{lstlisting}

\paragraph{jOOQ}
\begin{lstlisting}[frame=single,basicstyle=\small\ttfamily]
var employee = dslContext.select(DSL.asterisk())
    .from(EMPLOYEE)
    .where(EMPLOYEE.NAME.eq("Jack"))
    .fetchOne();
\end{lstlisting}

\clearpage
\section{Detection representations and output units}

Table~\ref{tab:detectionApproaches} contrasts the representations and output units of prior detection approaches with the class-labelled source spans evaluated here.

\begin{longtable}{>{\raggedright\arraybackslash}p{0.20\textwidth}>{\raggedright\arraybackslash}p{0.34\textwidth}>{\raggedright\arraybackslash}p{0.36\textwidth}}
    \caption{Detection representations and output units relevant to jOOQ source analysis. API denotes application programming interface, AST denotes abstract syntax tree, and LLM denotes large language model; the comparison shows that this study evaluates multiple class-labelled source spans rather than statement warnings or file classifications.}\label{tab:detectionApproaches}\\
    \hline
    \textbf{Approach} & \textbf{Representation and output} & \textbf{Relevance to this study} \\ \hline
    \endfirsthead
    \caption[]{Detection representations and output units (continued).}\\
    \hline
    \textbf{Approach} & \textbf{Representation and output} & \textbf{Relevance to this study} \\ \hline
    \endhead
    \hline
    \endfoot
    SQL extraction and parsing &
    SQL recovered from host-language source or query logs; parsed-statement warnings \citep{sharma2018smelly,nagy2017static,DBLP:conf/sigmod/DintyalaNA20}. &
    Does not directly expose occurrences represented by jOOQ domain-specific language calls or generated schema classes. \\
    API or AST rules &
    Host-language syntax and application programming interface (API) patterns or abstract syntax tree (AST) nodes; warnings attached to matched source nodes \citep{DBLP:journals/tse/MohaGDM10}. &
    Can localise encoded patterns but requires explicit coverage of jOOQ idioms and cross-file relationships. \\
    LLM source analysis &
    Source code and supplied context; prior work reports smell recognition or classification \citep{de2025effectiveness,blaauwendraad2025evaluating}. &
    Can return class-labelled spans, whose agreement and class-specific errors require empirical evaluation. \\
    This study &
    Java domain-specific language and generated schema code; multiple class-labelled source-line spans. &
    Evaluates the occurrence output later counted as detector flags across repositories. \\
\end{longtable}

\section{Project-disjoint split support}\label{sec:splitSupport}

Table~\ref{tab:splitSupport} shows the \num{21}/\num{20}/\num{20}-repository split and the sparse project support for several classes.

\begin{table}[H]
    \centering
    \small
    \caption{Occurrence and repository support in the \num{21}-repository training, \num{20}-repository validation, and \num{20}-repository test partitions. Each cell gives occurrences followed by repositories in parentheses; several classes occur in only one to five projects per partition.}
    \label{tab:splitSupport}
    \begin{tabular}{lrrrr}
        \hline
        \textbf{Class} & \textbf{Total} & \textbf{Training} & \textbf{Validation} & \textbf{Test} \\
        \hline
        Implicit Columns & 795 & 175 (20) & 350 (19) & 270 (20) \\
        ID Required & 259 & 55 (18) & 99 (15) & 105 (16) \\
        Keyless Entry & 122 & 30 (7) & 56 (1) & 36 (3) \\
        Fear of the Unknown & 95 & 23 (6) & 36 (8) & 36 (5) \\
        31 Flavors & 71 & 17 (6) & 15 (4) & 39 (1) \\
        Poor Man's Search Engine & 55 & 19 (5) & 15 (4) & 21 (3) \\
        Rounding Errors & 39 & 13 (4) & 10 (1) & 16 (2) \\
        \hline
    \end{tabular}
\end{table}

\section{Intersection-over-union sensitivity}\label{sec:iouSensitivity}

Table~\ref{tab:iouSensitivity} shows that pooled micro F1 varies by only \num{0.010} across the four tested localisation thresholds.

\begin{table}[H]
    \centering
    \caption{Sensitivity of pooled occurrence-level agreement to the line-span intersection-over-union (IoU) threshold. TP, FP, and FN denote true positives, false positives, and false negatives. Each row applies the same one-to-one matching rule to all \num{536} predictions and the original reference spans; micro F1 spans \num{0.863}--\num{0.873}.}
    \label{tab:iouSensitivity}
    \begin{tabular}{rrrrrrr}
        \hline
        \textbf{IoU} & \textbf{TP} & \textbf{FP} & \textbf{FN} & \textbf{Precision} & \textbf{Recall} & \textbf{F1} \\
        \hline
        0.25 & 462 & 74 & 61 & 0.862 & 0.883 & 0.873 \\
        0.50 & 460 & 76 & 63 & 0.858 & 0.880 & 0.869 \\
        0.75 & 457 & 79 & 66 & 0.853 & 0.874 & 0.863 \\
        1.00 & 457 & 79 & 66 & 0.853 & 0.874 & 0.863 \\
        \hline
    \end{tabular}
\end{table}

\bibliographystyle{../template/spbasic}
\bibliography{references}

\end{document}
